\documentclass[]{article}
\usepackage{lmodern}
\usepackage{amssymb,amsmath}
\usepackage{ifxetex,ifluatex}
\usepackage{fixltx2e} % provides \textsubscript
\ifnum 0\ifxetex 1\fi\ifluatex 1\fi=0 % if pdftex
  \usepackage[T1]{fontenc}
  \usepackage[utf8]{inputenc}
\else % if luatex or xelatex
  \ifxetex
    \usepackage{mathspec}
    \usepackage{xltxtra,xunicode}
  \else
    \usepackage{fontspec}
  \fi
  \defaultfontfeatures{Mapping=tex-text,Scale=MatchLowercase}
  \newcommand{\euro}{€}
\fi
% use upquote if available, for straight quotes in verbatim environments
\IfFileExists{upquote.sty}{\usepackage{upquote}}{}
% use microtype if available
\IfFileExists{microtype.sty}{\usepackage{microtype}}{}
\usepackage[margin=1in]{geometry}
\usepackage{color}
\usepackage{fancyvrb}
\newcommand{\VerbBar}{|}
\newcommand{\VERB}{\Verb[commandchars=\\\{\}]}
\DefineVerbatimEnvironment{Highlighting}{Verbatim}{commandchars=\\\{\}}
% Add ',fontsize=\small' for more characters per line
\newenvironment{Shaded}{}{}
\newcommand{\AlertTok}[1]{\textcolor[rgb]{1.00,0.00,0.00}{\textbf{#1}}}
\newcommand{\AnnotationTok}[1]{\textcolor[rgb]{0.38,0.63,0.69}{\textbf{\textit{#1}}}}
\newcommand{\AttributeTok}[1]{\textcolor[rgb]{0.49,0.56,0.16}{#1}}
\newcommand{\BaseNTok}[1]{\textcolor[rgb]{0.25,0.63,0.44}{#1}}
\newcommand{\BuiltInTok}[1]{\textcolor[rgb]{0.00,0.50,0.00}{#1}}
\newcommand{\CharTok}[1]{\textcolor[rgb]{0.25,0.44,0.63}{#1}}
\newcommand{\CommentTok}[1]{\textcolor[rgb]{0.38,0.63,0.69}{\textit{#1}}}
\newcommand{\CommentVarTok}[1]{\textcolor[rgb]{0.38,0.63,0.69}{\textbf{\textit{#1}}}}
\newcommand{\ConstantTok}[1]{\textcolor[rgb]{0.53,0.00,0.00}{#1}}
\newcommand{\ControlFlowTok}[1]{\textcolor[rgb]{0.00,0.44,0.13}{\textbf{#1}}}
\newcommand{\DataTypeTok}[1]{\textcolor[rgb]{0.56,0.13,0.00}{#1}}
\newcommand{\DecValTok}[1]{\textcolor[rgb]{0.25,0.63,0.44}{#1}}
\newcommand{\DocumentationTok}[1]{\textcolor[rgb]{0.73,0.13,0.13}{\textit{#1}}}
\newcommand{\ErrorTok}[1]{\textcolor[rgb]{1.00,0.00,0.00}{\textbf{#1}}}
\newcommand{\ExtensionTok}[1]{#1}
\newcommand{\FloatTok}[1]{\textcolor[rgb]{0.25,0.63,0.44}{#1}}
\newcommand{\FunctionTok}[1]{\textcolor[rgb]{0.02,0.16,0.49}{#1}}
\newcommand{\ImportTok}[1]{\textcolor[rgb]{0.00,0.50,0.00}{\textbf{#1}}}
\newcommand{\InformationTok}[1]{\textcolor[rgb]{0.38,0.63,0.69}{\textbf{\textit{#1}}}}
\newcommand{\KeywordTok}[1]{\textcolor[rgb]{0.00,0.44,0.13}{\textbf{#1}}}
\newcommand{\NormalTok}[1]{#1}
\newcommand{\OperatorTok}[1]{\textcolor[rgb]{0.40,0.40,0.40}{#1}}
\newcommand{\OtherTok}[1]{\textcolor[rgb]{0.00,0.44,0.13}{#1}}
\newcommand{\PreprocessorTok}[1]{\textcolor[rgb]{0.74,0.48,0.00}{#1}}
\newcommand{\RegionMarkerTok}[1]{#1}
\newcommand{\SpecialCharTok}[1]{\textcolor[rgb]{0.25,0.44,0.63}{#1}}
\newcommand{\SpecialStringTok}[1]{\textcolor[rgb]{0.73,0.40,0.53}{#1}}
\newcommand{\StringTok}[1]{\textcolor[rgb]{0.25,0.44,0.63}{#1}}
\newcommand{\VariableTok}[1]{\textcolor[rgb]{0.10,0.09,0.49}{#1}}
\newcommand{\VerbatimStringTok}[1]{\textcolor[rgb]{0.25,0.44,0.63}{#1}}
\newcommand{\WarningTok}[1]{\textcolor[rgb]{0.38,0.63,0.69}{\textbf{\textit{#1}}}}
\ifxetex
  \usepackage[setpagesize=false, % page size defined by xetex
              unicode=false, % unicode breaks when used with xetex
              xetex]{hyperref}
\else
  \usepackage[unicode=true]{hyperref}
\fi
\hypersetup{breaklinks=true,
            bookmarks=true,
            pdfauthor={},
            pdftitle={Sums and Subtypes and Unions},
            colorlinks=true,
            citecolor=blue,
            urlcolor=blue,
            linkcolor=magenta,
            pdfborder={0 0 0}}
\urlstyle{same}  % don't use monospace font for urls
% Make links footnotes instead of hotlinks:
\renewcommand{\href}[2]{#2\footnote{\url{#1}}}
\setlength{\parindent}{0pt}
\setlength{\parskip}{6pt plus 2pt minus 1pt}
\setlength{\emergencystretch}{3em}  % prevent overfull lines
\setcounter{secnumdepth}{0}

\title{Sums and Subtypes and Unions}

\begin{document}
\maketitle

\% Justin Le

\emph{Originally posted on
\textbf{\href{https://blog.jle.im/entry/sums-and-subtypes-and-unions.html}{in
Code}}.}

There's yet again been a bit of functional programming-adjacent twitter drama
recently, but it's actually sort of touched into some subtleties about sum types
that I am asked about (and think about) a lot nowadays. So, I'd like to take
this opportunity to talk a bit about the ``why'' and nature of sum types and how
to use them effectively, and how they contrast with other related concepts in
programming and software development and when even cases where sum types aren't
the best option.

\section{Sum Types at their Best}\label{sum-types-at-their-best}

The quintessential sum type that you just can't live without is \texttt{Maybe},
now adopted in a lot of languages as \texttt{Optional}:

\begin{Shaded}
\begin{Highlighting}[]
\KeywordTok{data} \DataTypeTok{Maybe}\NormalTok{ a }\OtherTok{=} \DataTypeTok{Nothing} \OperatorTok{|} \DataTypeTok{Just}\NormalTok{ a}
\end{Highlighting}
\end{Shaded}

If you have a value of type \texttt{Maybe\ Int}, it means that its valid values
are \texttt{Nothing}, \texttt{Just\ 0}, \texttt{Just\ 1}, etc.

This is also a good illustration to why we call it a ``sum'' type: if \texttt{a}
has \texttt{n} possible values, then \texttt{Maybe\ a} has \texttt{1\ +\ n}: we
add the single new value \texttt{Nothing} to it.

The ``benefit'' of the sum type is illustrated pretty clearly here too: every
time you \emph{use} a value of type \texttt{Maybe\ Int}, you are forced to
consider the fact that it could be \texttt{Nothing}:

\begin{Shaded}
\begin{Highlighting}[]
\OtherTok{showMaybeInt ::} \DataTypeTok{Maybe} \DataTypeTok{Int} \OtherTok{{-}\textgreater{}} \DataTypeTok{String}
\NormalTok{showMaybeInt }\OtherTok{=}\NormalTok{ \textbackslash{}}\KeywordTok{case}
  \DataTypeTok{Nothing} \OtherTok{{-}\textgreater{}} \StringTok{"There\textquotesingle{}s nothing here"}
  \DataTypeTok{Just}\NormalTok{ i }\OtherTok{{-}\textgreater{}} \StringTok{"Something is here: "} \OperatorTok{\textless{}\textgreater{}} \FunctionTok{show}\NormalTok{ i}
\end{Highlighting}
\end{Shaded}

That's because usually in sum type implementations, they are implemented in a
way that forces you to to handle each case exhaustively. Otherwise, sum types
are \emph{much} less useful.

At the most fundamental level this behaves like a compiler-enforced null check,
but built within the language in user-space (if it's implemented like a true sum
type) instead being compiler magic, ad-hoc syntax, or static analysis --- and
the fact that it can live in user-space is why it's been adopted so widely. At a
higher level, functional abstractions like Functor, Applicative, Monad,
Foldable, Traversable allow you to use a \texttt{Maybe\ a} like just a normal
\texttt{a} with the appropriate semantics, but that's
\href{https://blog.jle.im/entry/inside-my-world-ode-to-functor-and-monad.html}{a
topic for another time}.

This basic pattern can be extended to include more error information in your
\texttt{Nothing} branch, which is how you get the \texttt{Either\ e\ a} type in
the Haskell standard library, or the \texttt{Result\textless{}T,E\textgreater{}}
type in rust.

As a quick aside, remember that this example is a ``generic'' type (it has a
type parameter), but that's unrelated to the fact that it's a sum type. After
all, we could have a sum type \texttt{MaybeInt} that's either \texttt{NoInt} or
\texttt{JustInt\ 3}, \texttt{JustInt\ 4}, etc. I just wanted to clarify that the
``sum typeness'' (possible different distinguishable branches) is separate from
the ``parameterized typeenes''

Along different lines (not an error abstraction, not parameterized), we have the
common use case of defining syntax trees:

\begin{Shaded}
\begin{Highlighting}[]
\KeywordTok{data} \DataTypeTok{Expr} \OtherTok{=}
    \DataTypeTok{Lit} \DataTypeTok{Int}
  \OperatorTok{|} \DataTypeTok{Negate} \DataTypeTok{Expr}
  \OperatorTok{|} \DataTypeTok{Add} \DataTypeTok{Expr} \DataTypeTok{Expr}
  \OperatorTok{|} \DataTypeTok{Sub} \DataTypeTok{Expr} \DataTypeTok{Expr}
  \OperatorTok{|} \DataTypeTok{Mul} \DataTypeTok{Expr} \DataTypeTok{Expr}

\OtherTok{eval ::} \DataTypeTok{Expr} \OtherTok{{-}\textgreater{}} \DataTypeTok{Int}
\NormalTok{eval }\OtherTok{=}\NormalTok{ \textbackslash{}}\KeywordTok{case}
    \DataTypeTok{Lit}\NormalTok{ i }\OtherTok{{-}\textgreater{}}\NormalTok{ i}
    \DataTypeTok{Negate}\NormalTok{ x }\OtherTok{{-}\textgreater{}} \OperatorTok{{-}}\NormalTok{(eval x)}
    \DataTypeTok{Add}\NormalTok{ x y }\OtherTok{{-}\textgreater{}}\NormalTok{ eval x }\OperatorTok{+}\NormalTok{ eval y}
    \DataTypeTok{Sub}\NormalTok{ x y }\OtherTok{{-}\textgreater{}}\NormalTok{ eval x }\OperatorTok{{-}}\NormalTok{ eval y}
    \DataTypeTok{Mul}\NormalTok{ x y }\OtherTok{{-}\textgreater{}}\NormalTok{ eval x }\OperatorTok{*}\NormalTok{ eval y}

\OtherTok{pretty ::} \DataTypeTok{Expr} \OtherTok{{-}\textgreater{}} \DataTypeTok{String}
\NormalTok{pretty }\OtherTok{=}\NormalTok{ go }\DecValTok{0}
  \KeywordTok{where}
\OtherTok{    wrap ::} \DataTypeTok{Int} \OtherTok{{-}\textgreater{}} \DataTypeTok{Int} \OtherTok{{-}\textgreater{}} \DataTypeTok{String} \OtherTok{{-}\textgreater{}} \DataTypeTok{String}
\NormalTok{    wrap prio opPrec s}
      \OperatorTok{|}\NormalTok{ prio }\OperatorTok{\textgreater{}}\NormalTok{ opPrec }\OtherTok{=} \StringTok{"("} \OperatorTok{\textless{}\textgreater{}}\NormalTok{ s }\OperatorTok{\textless{}\textgreater{}} \StringTok{")"}
      \OperatorTok{|} \FunctionTok{otherwise} \OtherTok{=}\NormalTok{ s}
\NormalTok{    go prio }\OtherTok{=}\NormalTok{ \textbackslash{}}\KeywordTok{case}
        \DataTypeTok{Lit}\NormalTok{ i }\OtherTok{{-}\textgreater{}} \FunctionTok{show}\NormalTok{ i}
        \DataTypeTok{Negate}\NormalTok{ x }\OtherTok{{-}\textgreater{}}\NormalTok{ wrap prio }\DecValTok{2} \OperatorTok{$} \StringTok{"{-}"} \OperatorTok{\textless{}\textgreater{}}\NormalTok{ go }\DecValTok{2}\NormalTok{ x}
        \DataTypeTok{Add}\NormalTok{ x y }\OtherTok{{-}\textgreater{}}\NormalTok{ wrap prio }\DecValTok{0} \OperatorTok{$}\NormalTok{ go }\DecValTok{0}\NormalTok{ x }\OperatorTok{\textless{}\textgreater{}} \StringTok{" + "} \OperatorTok{\textless{}\textgreater{}}\NormalTok{ go }\DecValTok{1}\NormalTok{ y}
        \DataTypeTok{Sub}\NormalTok{ x y }\OtherTok{{-}\textgreater{}}\NormalTok{ wrap prio }\DecValTok{0} \OperatorTok{$}\NormalTok{ go }\DecValTok{0}\NormalTok{ x }\OperatorTok{\textless{}\textgreater{}} \StringTok{" {-} "} \OperatorTok{\textless{}\textgreater{}}\NormalTok{ go }\DecValTok{1}\NormalTok{ y}
        \DataTypeTok{Mul}\NormalTok{ x y }\OtherTok{{-}\textgreater{}}\NormalTok{ wrap prio }\DecValTok{1} \OperatorTok{$}\NormalTok{ go }\DecValTok{1}\NormalTok{ x }\OperatorTok{\textless{}\textgreater{}} \StringTok{" * "} \OperatorTok{\textless{}\textgreater{}}\NormalTok{ go }\DecValTok{2}\NormalTok{ y}

\OtherTok{main ::} \DataTypeTok{IO}\NormalTok{ ()}
\NormalTok{main }\OtherTok{=} \KeywordTok{do}
    \FunctionTok{putStrLn} \OperatorTok{$}\NormalTok{ pretty myExpr}
    \FunctionTok{print} \OperatorTok{$}\NormalTok{ eval myExpr}
  \KeywordTok{where}
\NormalTok{    myExpr }\OtherTok{=} \DataTypeTok{Mul}\NormalTok{ (}\DataTypeTok{Negate}\NormalTok{ (}\DataTypeTok{Add}\NormalTok{ (}\DataTypeTok{Lit} \DecValTok{4}\NormalTok{) (}\DataTypeTok{Lit} \DecValTok{5}\NormalTok{))) (}\DataTypeTok{Lit} \DecValTok{8}\NormalTok{)}
\end{Highlighting}
\end{Shaded}

\begin{verbatim}
-(4 + 5) * 8
-72
\end{verbatim}

Now, if we add a new command to the sum type, the compiler enforces us to handle
it.

\begin{Shaded}
\begin{Highlighting}[]
\KeywordTok{data} \DataTypeTok{Expr} \OtherTok{=}
    \DataTypeTok{Lit} \DataTypeTok{Int}
  \OperatorTok{|} \DataTypeTok{Negate} \DataTypeTok{Expr}
  \OperatorTok{|} \DataTypeTok{Add} \DataTypeTok{Expr} \DataTypeTok{Expr}
  \OperatorTok{|} \DataTypeTok{Sub} \DataTypeTok{Expr} \DataTypeTok{Expr}
  \OperatorTok{|} \DataTypeTok{Mul} \DataTypeTok{Expr} \DataTypeTok{Expr}
  \OperatorTok{|} \DataTypeTok{Abs} \DataTypeTok{Expr}

\OtherTok{eval ::} \DataTypeTok{Expr} \OtherTok{{-}\textgreater{}} \DataTypeTok{Int}
\NormalTok{eval }\OtherTok{=}\NormalTok{ \textbackslash{}}\KeywordTok{case}
    \DataTypeTok{Lit}\NormalTok{ i }\OtherTok{{-}\textgreater{}}\NormalTok{ i}
    \DataTypeTok{Negate}\NormalTok{ x }\OtherTok{{-}\textgreater{}} \OperatorTok{{-}}\NormalTok{(eval x)}
    \DataTypeTok{Add}\NormalTok{ x y }\OtherTok{{-}\textgreater{}}\NormalTok{ eval x }\OperatorTok{+}\NormalTok{ eval y}
    \DataTypeTok{Sub}\NormalTok{ x y }\OtherTok{{-}\textgreater{}}\NormalTok{ eval x }\OperatorTok{{-}}\NormalTok{ eval y}
    \DataTypeTok{Mul}\NormalTok{ x y }\OtherTok{{-}\textgreater{}}\NormalTok{ eval x }\OperatorTok{*}\NormalTok{ eval y}
    \DataTypeTok{Abs}\NormalTok{ x }\OtherTok{{-}\textgreater{}} \FunctionTok{abs}\NormalTok{ (eval x)}

\OtherTok{pretty ::} \DataTypeTok{Expr} \OtherTok{{-}\textgreater{}} \DataTypeTok{String}
\NormalTok{pretty }\OtherTok{=}\NormalTok{ go }\DecValTok{0}
  \KeywordTok{where}
\OtherTok{    wrap ::} \DataTypeTok{Int} \OtherTok{{-}\textgreater{}} \DataTypeTok{Int} \OtherTok{{-}\textgreater{}} \DataTypeTok{String} \OtherTok{{-}\textgreater{}} \DataTypeTok{String}
\NormalTok{    wrap prio opPrec s}
      \OperatorTok{|}\NormalTok{ prio }\OperatorTok{\textgreater{}}\NormalTok{ opPrec }\OtherTok{=} \StringTok{"("} \OperatorTok{\textless{}\textgreater{}}\NormalTok{ s }\OperatorTok{\textless{}\textgreater{}} \StringTok{")"}
      \OperatorTok{|} \FunctionTok{otherwise} \OtherTok{=}\NormalTok{ s}
\NormalTok{    go prio }\OtherTok{=}\NormalTok{ \textbackslash{}}\KeywordTok{case}
        \DataTypeTok{Lit}\NormalTok{ i }\OtherTok{{-}\textgreater{}} \FunctionTok{show}\NormalTok{ i}
        \DataTypeTok{Negate}\NormalTok{ x }\OtherTok{{-}\textgreater{}}\NormalTok{ wrap prio }\DecValTok{2} \OperatorTok{$} \StringTok{"{-}"} \OperatorTok{\textless{}\textgreater{}}\NormalTok{ go }\DecValTok{2}\NormalTok{ x}
        \DataTypeTok{Add}\NormalTok{ x y }\OtherTok{{-}\textgreater{}}\NormalTok{ wrap prio }\DecValTok{0} \OperatorTok{$}\NormalTok{ go }\DecValTok{0}\NormalTok{ x }\OperatorTok{\textless{}\textgreater{}} \StringTok{" + "} \OperatorTok{\textless{}\textgreater{}}\NormalTok{ go }\DecValTok{1}\NormalTok{ y}
        \DataTypeTok{Sub}\NormalTok{ x y }\OtherTok{{-}\textgreater{}}\NormalTok{ wrap prio }\DecValTok{0} \OperatorTok{$}\NormalTok{ go }\DecValTok{0}\NormalTok{ x }\OperatorTok{\textless{}\textgreater{}} \StringTok{" {-} "} \OperatorTok{\textless{}\textgreater{}}\NormalTok{ go }\DecValTok{1}\NormalTok{ y}
        \DataTypeTok{Mul}\NormalTok{ x y }\OtherTok{{-}\textgreater{}}\NormalTok{ wrap prio }\DecValTok{1} \OperatorTok{$}\NormalTok{ go }\DecValTok{1}\NormalTok{ x }\OperatorTok{\textless{}\textgreater{}} \StringTok{" * "} \OperatorTok{\textless{}\textgreater{}}\NormalTok{ go }\DecValTok{2}\NormalTok{ y}
        \DataTypeTok{Abs}\NormalTok{ x }\OtherTok{{-}\textgreater{}}\NormalTok{ wrap prio }\DecValTok{2} \OperatorTok{$} \StringTok{"|"} \OperatorTok{\textless{}\textgreater{}}\NormalTok{ go }\DecValTok{0}\NormalTok{ x }\OperatorTok{\textless{}\textgreater{}} \StringTok{"|"}
\end{Highlighting}
\end{Shaded}

Another example where things shine are as clearly-fined APIs between processes.
For example, we can imagine a ``command'' type that sends different types of
commands with different payloads. This can be interpreted as perhaps the result
of parsing command line arguments or the message in some communication protocol.

For example, you could have a protocol that launches and controls processes:

\begin{Shaded}
\begin{Highlighting}[]
\KeywordTok{data} \DataTypeTok{Command}\NormalTok{ a }\OtherTok{=}
    \DataTypeTok{Launch} \DataTypeTok{String}\NormalTok{ (}\DataTypeTok{Int} \OtherTok{{-}\textgreater{}}\NormalTok{ a)    }\CommentTok{{-}{-} \^{} takes a name, returns a process ID}
  \OperatorTok{|} \DataTypeTok{Stop} \DataTypeTok{Int}\NormalTok{ (}\DataTypeTok{Bool} \OtherTok{{-}\textgreater{}}\NormalTok{ a)        }\CommentTok{{-}{-} \^{} takes a process ID, returns success/failure}
\end{Highlighting}
\end{Shaded}

This ADT is written in the ``interpreter'' pattern, where any arguments not
involving \texttt{a} are the command payload, any \texttt{X\ -\textgreater{}\ a}
represent that the command could respond with \texttt{X}.

Let's write a sample interpreter backing the state in an IntMap in an IORef:

\begin{Shaded}
\begin{Highlighting}[]
\KeywordTok{import} \KeywordTok{qualified} \DataTypeTok{Data.IntMap} \KeywordTok{as} \DataTypeTok{IM}
\KeywordTok{import} \DataTypeTok{Data.IntMap}\NormalTok{ (}\DataTypeTok{IntMap}\NormalTok{)}

\OtherTok{runCommand ::} \DataTypeTok{IORef}\NormalTok{ (}\DataTypeTok{IntMap} \DataTypeTok{String}\NormalTok{) }\OtherTok{{-}\textgreater{}} \DataTypeTok{Command}\NormalTok{ a }\OtherTok{{-}\textgreater{}} \DataTypeTok{IO}\NormalTok{ a}
\NormalTok{runCommand ref }\OtherTok{=}\NormalTok{ \textbackslash{}}\KeywordTok{case}
    \DataTypeTok{Launch}\NormalTok{ newName next }\OtherTok{{-}\textgreater{}} \KeywordTok{do}
\NormalTok{        currMap }\OtherTok{\textless{}{-}}\NormalTok{ readIORef ref}
        \KeywordTok{let}\NormalTok{ newId }\OtherTok{=} \KeywordTok{case}\NormalTok{ IM.lookupMax currMap }\KeywordTok{of}
              \DataTypeTok{Nothing} \OtherTok{{-}\textgreater{}} \DecValTok{0}
              \DataTypeTok{Just}\NormalTok{ (i, \_) }\OtherTok{{-}\textgreater{}}\NormalTok{ i }\OperatorTok{+} \DecValTok{1}
\NormalTok{        modifyIORef ref }\OperatorTok{$}\NormalTok{ IM.insert newId newName}
        \FunctionTok{pure}\NormalTok{ (next newId)}
    \DataTypeTok{Stop}\NormalTok{ procId next }\OtherTok{{-}\textgreater{}} \KeywordTok{do}
\NormalTok{        existed }\OtherTok{\textless{}{-}}\NormalTok{ IM.member procId }\OperatorTok{\textless{}$\textgreater{}}\NormalTok{ readIORef ref}
\NormalTok{        modifyIORef ref }\OperatorTok{$}\NormalTok{ IM.delete procId}
        \FunctionTok{pure}\NormalTok{ (next existed)}

\OtherTok{main ::} \DataTypeTok{IO}\NormalTok{ ()}
\NormalTok{main }\OtherTok{=} \KeywordTok{do}
\NormalTok{    ref }\OtherTok{\textless{}{-}}\NormalTok{ newIORef IM.empty}
\NormalTok{    aliceId }\OtherTok{\textless{}{-}}\NormalTok{ runCommand ref }\OperatorTok{$} \DataTypeTok{Launch} \StringTok{"alice"} \FunctionTok{id}
    \FunctionTok{putStrLn} \OperatorTok{$} \StringTok{"Launched alice with ID "} \OperatorTok{\textless{}\textgreater{}} \FunctionTok{show}\NormalTok{ aliceId}
\NormalTok{    bobId }\OtherTok{\textless{}{-}}\NormalTok{ runCommand ref }\OperatorTok{$} \DataTypeTok{Launch} \StringTok{"bob"} \FunctionTok{id}
    \FunctionTok{putStrLn} \OperatorTok{$} \StringTok{"Launched bob with ID "} \OperatorTok{\textless{}\textgreater{}} \FunctionTok{show}\NormalTok{ bobId}
\NormalTok{    success }\OtherTok{\textless{}{-}}\NormalTok{ runCommand ref }\OperatorTok{$} \DataTypeTok{Stop}\NormalTok{ aliceId }\FunctionTok{id}
    \FunctionTok{putStrLn} \OperatorTok{$}
      \KeywordTok{if}\NormalTok{ success}
        \KeywordTok{then} \StringTok{"alice succesfully stopped"}
        \KeywordTok{else} \StringTok{"alice unsuccesfully stopped"}
    \FunctionTok{print} \OperatorTok{=\textless{}\textless{}}\NormalTok{ readIORef ref}
\end{Highlighting}
\end{Shaded}

\begin{verbatim}
Launched alice with ID 0
Launched bob with ID 0
alice succesfully stopped
fromList [(1, "bob")]
\end{verbatim}

In a real application you might interpret this command in \texttt{IO}, but for
demonstration's sake, here's a way of ``handling'' this command purely using an
\texttt{IntMap} of id's to names to back the state:

Let's add a command to ``query'' a process id for its current status:

\begin{Shaded}
\begin{Highlighting}[]
\KeywordTok{data} \DataTypeTok{Command}\NormalTok{ a }\OtherTok{=}
    \DataTypeTok{Launch} \DataTypeTok{String}\NormalTok{ (}\DataTypeTok{Int} \OtherTok{{-}\textgreater{}}\NormalTok{ a)    }\CommentTok{{-}{-} \^{} takes a name, returns a process ID}
  \OperatorTok{|} \DataTypeTok{Stop} \DataTypeTok{Int}\NormalTok{ (}\DataTypeTok{Bool} \OtherTok{{-}\textgreater{}}\NormalTok{ a)        }\CommentTok{{-}{-} \^{} takes a process ID, returns success/failure}
  \OperatorTok{|} \DataTypeTok{Query} \DataTypeTok{Int}\NormalTok{ (}\DataTypeTok{String} \OtherTok{{-}\textgreater{}}\NormalTok{ a)     }\CommentTok{{-}{-} \^{} takes a process ID, returns a status message}

\OtherTok{runCommand ::} \DataTypeTok{IORef}\NormalTok{ (}\DataTypeTok{IntMap} \DataTypeTok{String}\NormalTok{) }\OtherTok{{-}\textgreater{}} \DataTypeTok{Command}\NormalTok{ a }\OtherTok{{-}\textgreater{}} \DataTypeTok{IO}\NormalTok{ a}
\NormalTok{runCommand ref }\OtherTok{=}\NormalTok{ \textbackslash{}}\KeywordTok{case}
    \DataTypeTok{Launch}\NormalTok{ newName next }\OtherTok{{-}\textgreater{}} \KeywordTok{do}
\NormalTok{        currMap }\OtherTok{\textless{}{-}}\NormalTok{ readIORef ref}
        \KeywordTok{let}\NormalTok{ newId }\OtherTok{=} \KeywordTok{case}\NormalTok{ IM.lookupMax currMap }\KeywordTok{of}
              \DataTypeTok{Nothing} \OtherTok{{-}\textgreater{}} \DecValTok{0}
              \DataTypeTok{Just}\NormalTok{ (i, \_) }\OtherTok{{-}\textgreater{}}\NormalTok{ i }\OperatorTok{+} \DecValTok{1}
\NormalTok{        modifyIORef ref }\OperatorTok{$}\NormalTok{ IM.insert newId newName}
        \FunctionTok{pure}\NormalTok{ (next newId)}
    \DataTypeTok{Stop}\NormalTok{ procId next }\OtherTok{{-}\textgreater{}} \KeywordTok{do}
\NormalTok{        existed }\OtherTok{\textless{}{-}}\NormalTok{ IM.member procId }\OperatorTok{\textless{}$\textgreater{}}\NormalTok{ readIORef ref}
\NormalTok{        modify }\OperatorTok{$}\NormalTok{ IM.delete procId}
        \FunctionTok{pure}\NormalTok{ (next existed)}
    \DataTypeTok{Query}\NormalTok{ procId next }\OtherTok{{-}\textgreater{}} \KeywordTok{do}
\NormalTok{        procName }\OtherTok{\textless{}{-}}\NormalTok{ IM.lookup procId }\OperatorTok{\textless{}$\textgreater{}}\NormalTok{ readIORef ref}
        \FunctionTok{pure} \KeywordTok{case}\NormalTok{ procName }\KeywordTok{of}
          \DataTypeTok{Nothing} \OtherTok{{-}\textgreater{}} \StringTok{"This process doesn\textquotesingle{}t exist, silly."}
          \DataTypeTok{Just}\NormalTok{ n }\OtherTok{{-}\textgreater{}} \StringTok{"Process "} \OperatorTok{\textless{}\textgreater{}}\NormalTok{ n }\OperatorTok{\textless{}\textgreater{}} \StringTok{" chugging along..."}
\end{Highlighting}
\end{Shaded}

\subsection{Relationship with Unions}\label{relationship-with-unions}

To clarify a common confusion, sum types could be described as ``tagged
unions'': you have a tag to indicate which branch you are on (which can be
case-matched on), and then the rest of your data is conditionally present.

In many languages this can be implemented \emph{literally} as a struct with a
tag and a union of data.

Remember, it's not \emph{exactly} a union, because, ie, consider a type like:

\begin{Shaded}
\begin{Highlighting}[]
\KeywordTok{data} \DataTypeTok{Entity} \OtherTok{=} \DataTypeTok{User} \DataTypeTok{Int} \OperatorTok{|} \DataTypeTok{Post} \DataTypeTok{Int}
\end{Highlighting}
\end{Shaded}

This data could represent a user at a user id, or a post at a post id. If we
considered it purely as a union of \texttt{Int} and \texttt{Int}:

\begin{Shaded}
\begin{Highlighting}[]
\KeywordTok{union}\NormalTok{ Entity }\OperatorTok{\{}
    \DataTypeTok{int}\NormalTok{ user\_id}\OperatorTok{;}
    \DataTypeTok{int}\NormalTok{ post\_id}\OperatorTok{;}
\OperatorTok{\};}
\end{Highlighting}
\end{Shaded}

we'd lose the ability to branch on whether or not we have a user or an int. If
we have the tagged union, we recover the original tagged union semantics:

\begin{Shaded}
\begin{Highlighting}[]
\KeywordTok{struct}\NormalTok{ Entity }\OperatorTok{\{}
    \DataTypeTok{bool}\NormalTok{ is\_user}\OperatorTok{;}
    \KeywordTok{union} \OperatorTok{\{}
        \DataTypeTok{int}\NormalTok{ user\_id}\OperatorTok{;}
        \DataTypeTok{int}\NormalTok{ post\_id}\OperatorTok{;}
    \OperatorTok{\}}\NormalTok{ payload}\OperatorTok{;}
\OperatorTok{\};}
\end{Highlighting}
\end{Shaded}

Of course, you still need an abstract interface like the visitor pattern to
actually be able to use this as a sum type with guarantees that you handle every
branch, but that's a story for another day. Alternatively, if your language
supports dynamic dispatch nicely, that's another underlying implementation that
would work to back a higher-level visitor pattern interface.

\section{Subtypes Solve a Different
Problem}\label{subtypes-solve-a-different-problem}

Now, sum types aren't exactly a part of common programming education curriculum,
but \emph{subtypes} and \emph{supertypes} definitely were drilled into every CS
student's brain and waking nightmares from their first year.

Informally (a la Liskov), \texttt{B} is a subtype of \texttt{A} (and \texttt{A}
is a supertype of \texttt{B}) if anywhere that expects an \texttt{A}, you could
also provide a \texttt{B}.

In normal object-oriented programming, this often shows up in early lessons as
\texttt{Cat} and \texttt{Dog} being subclasses of an \texttt{Animal} class, or
\texttt{Square} and \texttt{Circle} being subclasses of a \texttt{Shape} class.

When people first learn about sum types, there is a tendency to understand them
as similar to subtyping. This is unfortunately understandable, since a lot of
introductions to sum types often start with something like

\begin{Shaded}
\begin{Highlighting}[]
\CommentTok{{-}{-} | Bad Sum Type Example!}
\KeywordTok{data} \DataTypeTok{Shape} \OtherTok{=} \DataTypeTok{Circle} \DataTypeTok{Double} \OperatorTok{|} \DataTypeTok{Rectangle} \DataTypeTok{Double} \DataTypeTok{Double}
\end{Highlighting}
\end{Shaded}

While there are situations where this might be a good sum type (ie, for an API
specification or a state machine), on face-value this is a bad example on the
sum types vs.~subtyping distinction.

You might notice the essential ``tension'' of the sum type: you declare all of
your options up-front, the functions that consume your value are open and
declared ad-hoc. And, if you add new options, all of the consuming functions
must be adjusted.

So, \emph{subtypes} (and supertypes) are more effective when they lean into the
opposite end: the universe of possible options are open and declared ad-hoc, but
the \emph{consuming functions} are closed. And, if you add new functions, all of
the members must be adjusted.

\subsection{Extensible Abstractions}\label{extensible-abstractions}

In Haskell, subtyping is implemented in terms of parametric polymorphism.

For example, consider a common API for json serialization. You might have a
bunch of things that can be ``empty'':

\begin{Shaded}
\begin{Highlighting}[]
\OtherTok{emptyString ::} \DataTypeTok{String}
\NormalTok{emptyString }\OtherTok{=} \StringTok{""}

\OtherTok{emptyList ::}\NormalTok{ [a]}
\NormalTok{emptyList }\OtherTok{=}\NormalTok{ []}

\OtherTok{emptyMaybe ::} \DataTypeTok{Maybe}\NormalTok{ a}
\NormalTok{emptyMaybe }\OtherTok{=} \DataTypeTok{Nothing}
\end{Highlighting}
\end{Shaded}

But, consider the \texttt{Monoid} typeclass, which gives you:

\begin{Shaded}
\begin{Highlighting}[]
\FunctionTok{mempty}\OtherTok{ ::} \DataTypeTok{Monoid}\NormalTok{ a }\OtherTok{=\textgreater{}}\NormalTok{ a}
\end{Highlighting}
\end{Shaded}

Now, any function that expects a \texttt{String}, \texttt{{[}Int{]}},
\texttt{Maybe\ Double}, etc., you can \emph{also} pass in \texttt{mempty}. So,
this means that the type \texttt{forall\ a.\ Monoid\ a\ =\textgreater{}\ a} is a
\emph{subtype} of \texttt{String}, \texttt{{[}Bool{]}}, etc.

Because of how typeclasses work, we can create a new supertype of
\texttt{forall\ a.\ Monoid\ a\ =\textgreater{}\ a}:

\begin{Shaded}
\begin{Highlighting}[]
\KeywordTok{newtype} \DataTypeTok{SumInt} \OtherTok{=} \DataTypeTok{Sum} \DataTypeTok{Int}

\KeywordTok{instance} \DataTypeTok{Semigroup} \DataTypeTok{SumInt} \KeywordTok{where}
    \DataTypeTok{SumInt}\NormalTok{ x }\OperatorTok{\textless{}\textgreater{}} \DataTypeTok{SumInt}\NormalTok{ y }\OtherTok{=} \DataTypeTok{SumInt}\NormalTok{ (x }\OperatorTok{+}\NormalTok{ y)}

\KeywordTok{instance} \DataTypeTok{Monoid} \DataTypeTok{SumInt} \KeywordTok{where}
    \FunctionTok{mempty} \OtherTok{=} \DataTypeTok{SumInt} \DecValTok{0}
\end{Highlighting}
\end{Shaded}

Another common example is to put the typeclass constraint on a type in the
negative (contravariant) part of the type. For example, you could have multiple
functions that serialize into JSON:

\begin{Shaded}
\begin{Highlighting}[]
\OtherTok{fooToJson ::} \DataTypeTok{Foo} \OtherTok{{-}\textgreater{}} \DataTypeTok{Value}
\OtherTok{barToJson ::} \DataTypeTok{Bar} \OtherTok{{-}\textgreater{}} \DataTypeTok{Value}
\OtherTok{bazToJson ::} \DataTypeTok{Baz} \OtherTok{{-}\textgreater{}} \DataTypeTok{Value}
\end{Highlighting}
\end{Shaded}

Through typeclasses, you can create:

\begin{Shaded}
\begin{Highlighting}[]
\OtherTok{toJSON ::} \DataTypeTok{ToJSON}\NormalTok{ a }\OtherTok{=\textgreater{}}\NormalTok{ a }\OtherTok{{-}\textgreater{}} \DataTypeTok{Value}
\end{Highlighting}
\end{Shaded}

The type of
\texttt{toJSON\ ::\ forall\ a.\ JSON\ a\ =\textgreater{}\ a\ -\textgreater{}\ Value}
is a subtype of \texttt{Foo\ -\textgreater{}\ Value},
\texttt{Bar\ -\textgreater{}\ Value}, and \texttt{Baz\ -\textgreater{}\ Value},
because everywhere you would \emph{want} a \texttt{Foo\ -\textgreater{}\ Value},
you could give \texttt{toJSON} instead. Every time you \emph{want} to serialize
a \texttt{Foo}, you could use \texttt{toJSON}.

This usage works well, as it gives you an extensible abstraction to design code
around. When you write code polymorphic over \texttt{Monoid\ a}, it forces you
to reason about your values with respect to only the aspects relating to
monoidness. If you write code polymorphic over \texttt{Num\ a}, it forces you to
reason about your values only with respect to how they can be added, subtracted,
negated, or multiplied, instead of having to worry about things like their
machine representation.

All in all, this is convenient because you can create \emph{even more
supertypes} of
\texttt{forall\ a.\ ToJSON\ a\ =\textgreater{}\ a\ -\textgreater{}\ Value}
easily, just by defining a new typeclass instance. So, if you want \emph{more}
things you can use \texttt{toJSON} on (or more things you can provide with
\texttt{mempty}), you just need to define the typeclass instance.

If you ever added a new method to \texttt{Monoid} or \texttt{ToJSON}, this is a
huge breaking change for all instances. But adding a new instance is easy!

\subsection{Containers of an Interface}\label{containers-of-an-interface}

A common abstraction in languages that rely heavily on subtyping is a
heterogeneous ``container'' (like a list) of things of different type where all
you know about each value is that they instantiate some interface.

For example, in Haskell, you could have a list of things that are
\texttt{Show}-able:

\begin{Shaded}
\begin{Highlighting}[]
\KeywordTok{data} \DataTypeTok{Showable} \OtherTok{=} \KeywordTok{forall}\NormalTok{ a}\OperatorTok{.} \DataTypeTok{Show}\NormalTok{ a }\OtherTok{=\textgreater{}} \DataTypeTok{Showable}\NormalTok{ a}
\end{Highlighting}
\end{Shaded}

So you could have a function like:

\begin{Shaded}
\begin{Highlighting}[]
\OtherTok{printAll ::}\NormalTok{ [}\DataTypeTok{Showable}\NormalTok{] }\OtherTok{{-}\textgreater{}} \DataTypeTok{IO}\NormalTok{ ()}
\NormalTok{printAll }\OtherTok{=}\NormalTok{ traverse\_ \textbackslash{}(}\DataTypeTok{Showable}\NormalTok{ x) }\OtherTok{{-}\textgreater{}} \FunctionTok{print}\NormalTok{ x}
\end{Highlighting}
\end{Shaded}

\begin{Shaded}
\begin{Highlighting}[]
\NormalTok{ghci}\OperatorTok{\textgreater{}}\NormalTok{ printAll [}\DataTypeTok{Showable} \DecValTok{3}\NormalTok{, }\DataTypeTok{Showable} \DataTypeTok{True}\NormalTok{, }\DataTypeTok{Showable} \StringTok{"hello"}\NormalTok{]}
\DecValTok{3}
\DataTypeTok{True}
\StringTok{"hello"}
\end{Highlighting}
\end{Shaded}

This \emph{specific} example is pretty silly because \texttt{{[}Showable{]}} is
just the same as \texttt{{[}String{]}}:

\begin{Shaded}
\begin{Highlighting}[]
\NormalTok{ghci}\OperatorTok{\textgreater{}}\NormalTok{ traverse\_ }\FunctionTok{putStrLn}\NormalTok{ [}\FunctionTok{show} \DecValTok{3}\NormalTok{, }\FunctionTok{show} \DataTypeTok{True}\NormalTok{, }\FunctionTok{show} \StringTok{"hello"}\NormalTok{]}
\DecValTok{3}
\DataTypeTok{True}
\StringTok{"hello"}
\end{Highlighting}
\end{Shaded}

However, this ``widget pattern'' does become more useful when your typeclass is
more complicated, like the \texttt{Layout} class in
\emph{\href{https://hackage.haskell.org/package/xmonad}{xmonad}}, where it can
be used to pass a container of ``widgets'' that are to be rendered, where the
widgets can all be different types but all implement a specific typeclass.
Usually you can also get away with a container of handler functions instead (so
you don't have to use existential types), but using typeclasses in this case
gives you guarantees of canonicity (one instance per nominal type) and some
convenient wrapping/unwraping and interactions between subclasses.

I do mention in
\href{https://blog.jle.im/entry/levels-of-type-safety-haskell-lists.html}{a blog
post about different types of existential lists}, however, that this ``container
of instances'' type is much less useful in Haskell than in other languages for
two main reasons. First, because Haskell gives you a whole wealth of
functionality to operate over homogeneous parameters (like \texttt{{[}a{]}},
where all items have the same type) that jumping to heterogeneous lists gives up
so much. Second, that ``widget container'' patterns in other languages often
resort to runtime reflection of time information, which means in practice you'd
have to add an extra \texttt{Typeable} constraint to your existentials to use
the same way you'd use them in OOP languages.

\begin{center}\rule{0.5\linewidth}{0.5pt}\end{center}

\textbf{Aside}

Let's briefly take a moment to talk about how typeclass hierarchies give us
subtle subtype/supertype relationships.

Let's look at the classic \texttt{Num} and \texttt{Fractional}:

\begin{Shaded}
\begin{Highlighting}[]
\KeywordTok{class} \DataTypeTok{Num}\NormalTok{ a}

\KeywordTok{class} \DataTypeTok{Num}\NormalTok{ a }\OtherTok{=\textgreater{}} \DataTypeTok{Fractional}\NormalTok{ a}
\end{Highlighting}
\end{Shaded}

\texttt{Num} is a \emph{superclass} of \texttt{Fractional}, and
\texttt{Fractional} is a \emph{subclass} of \texttt{Num}. Everywhere a
\texttt{Num} constraint is required, you can provide a \texttt{Fractional}
constraint to do the same thing.

However, in these two types:

\begin{Shaded}
\begin{Highlighting}[]
\DataTypeTok{Num}\NormalTok{ a }\OtherTok{=\textgreater{}}\NormalTok{ a}
\DataTypeTok{Fractional}\NormalTok{ a }\OtherTok{=\textgreater{}}\NormalTok{ a}
\end{Highlighting}
\end{Shaded}

\texttt{forall\ a.\ Num\ a\ =\textgreater{}\ a} is actually a \emph{subclass} of
\texttt{forall\ a.\ Fractional\ a\ =\textgreater{}\ a}! That's because if you
need a \texttt{forall\ a.\ Fractional\ a\ =\textgreater{}\ a}, you can provide a
\texttt{forall\ a.\ Num\ a\ =\textgreater{}\ a} instead. In fact, let's look at
three levels: \texttt{Double},
\texttt{forall\ a.\ Fractional\ a\ =\textgreater{}\ a}, and
\texttt{forall\ a.\ Num\ a\ =\textgreater{}\ a}.

\begin{Shaded}
\begin{Highlighting}[]
\CommentTok{{-}{-} can be used as \textasciigrave{}Double\textasciigrave{}}
\FloatTok{1.0}\OtherTok{ ::} \DataTypeTok{Double}
\FloatTok{1.0}\OtherTok{ ::} \DataTypeTok{Fractional}\NormalTok{ a }\OtherTok{=\textgreater{}}\NormalTok{ a}
\DecValTok{1}\OtherTok{ ::} \DataTypeTok{Num}\NormalTok{ a }\OtherTok{=\textgreater{}}\NormalTok{ a}

\CommentTok{{-}{-} can be used as \textasciigrave{}forall a. Fractional a =\textgreater{} a\textasciigrave{}}
\FloatTok{1.0}\OtherTok{ ::} \DataTypeTok{Fractional}\NormalTok{ a }\OtherTok{=\textgreater{}}\NormalTok{ a}
\DecValTok{1}\OtherTok{ ::} \DataTypeTok{Num}\NormalTok{ a }\OtherTok{=\textgreater{}}\NormalTok{ a}

\CommentTok{{-}{-} can be used as \textasciigrave{}forall a. Num a =\textgreater{} a\textasciigrave{}}
\DecValTok{1}\OtherTok{ ::} \DataTypeTok{Num}\NormalTok{ a }\OtherTok{=\textgreater{}}\NormalTok{ a}
\end{Highlighting}
\end{Shaded}

So, \texttt{Double} is a supertype of \texttt{Fractional\ a\ =\textgreater{}\ a}
is a supertype of \texttt{Num\ a\ =\textgreater{}\ a}.

The general idea here is that the more super- you go, the more you ``know''
about the actual term you are creating. So, with
\texttt{Num\ a\ =\textgreater{}\ a}, you know the \emph{least} (and, you have
the most possible actual terms because there are more instances of \texttt{Num}
than of \texttt{Fractional}). And, with \texttt{Double}, you know the
\emph{most}: you even know its machine representation!

So, \texttt{Num} is a superclass of \texttt{Fractional} but
\texttt{forall\ a.\ Num\ a\ =\textgreater{}\ a} is a subclass of
\texttt{forall\ a.\ Fractional\ a\ =\textgreater{}\ a}. This actually follows
the typical rules of subtyping: if something appears on the ``left'' of an arrow
(\texttt{=\textgreater{}} in this case), it gets flipped from sub- to super-. We
often call the left side a ``negative'' (contravariant) position and the right
side a ``positive'' position, because a negative of a negative (the left side of
a left size, like \texttt{a} in
\texttt{(a\ -\textgreater{}\ b)\ -\textgreater{}\ c}) is a positive.

Also note that our ``existential wrappers'':

\begin{Shaded}
\begin{Highlighting}[]
\KeywordTok{data} \DataTypeTok{SomeNum} \OtherTok{=} \KeywordTok{forall}\NormalTok{ a}\OperatorTok{.} \DataTypeTok{Num}\NormalTok{ a }\OtherTok{=\textgreater{}} \DataTypeTok{SomeFractional}\NormalTok{ a}
\KeywordTok{data} \DataTypeTok{SomeFractional} \OtherTok{=} \KeywordTok{forall}\NormalTok{ a}\OperatorTok{.} \DataTypeTok{Fractional}\NormalTok{ a }\OtherTok{=\textgreater{}} \DataTypeTok{SomeFractional}\NormalTok{ a}
\end{Highlighting}
\end{Shaded}

can be CPS-transformed to their equivalent types:

\begin{Shaded}
\begin{Highlighting}[]
\KeywordTok{type} \DataTypeTok{SomeNum} \OtherTok{=} \KeywordTok{forall}\NormalTok{ r}\OperatorTok{.}\NormalTok{ (}\KeywordTok{forall}\NormalTok{ a}\OperatorTok{.} \DataTypeTok{Num}\NormalTok{ a }\OtherTok{=\textgreater{}}\NormalTok{ a }\OtherTok{{-}\textgreater{}}\NormalTok{ r) }\OtherTok{{-}\textgreater{}}\NormalTok{ r}
\KeywordTok{type} \DataTypeTok{SomeFractional} \OtherTok{=} \KeywordTok{forall}\NormalTok{ r}\OperatorTok{.}\NormalTok{ (}\KeywordTok{forall}\NormalTok{ a}\OperatorTok{.} \DataTypeTok{Fractional}\NormalTok{ a }\OtherTok{=\textgreater{}}\NormalTok{ a }\OtherTok{{-}\textgreater{}}\NormalTok{ r) }\OtherTok{{-}\textgreater{}}\NormalTok{ r}
\end{Highlighting}
\end{Shaded}

And in those cases, \texttt{Num} and \texttt{Fractional} again appear in the
covariant (positive) position, since they're the negative of negative. So, this
aligns with our intuition that \texttt{SomeFractional} is a subtype of
\texttt{SomeNum}.

\begin{center}\rule{0.5\linewidth}{0.5pt}\end{center}

\section{The Expression Problem}\label{the-expression-problem}

This tension that I described earlier is often called
\href{https://en.wikipedia.org/wiki/Expression_problem}{The Expression Problem},
and is a tension that is inherent to a lot of aspects of language and
abstraction design. However, in the context laid out in this post, it serves as
a good general guide to decide what pattern to go down:

\begin{itemize}
\tightlist
\item
  If you expect a canonical set of ``inhabitants'' and an open set of
  ``operations'', sum types can suit that end of the spectrum well.
\item
  If you expect a canonical set of ``operations'' and an open set of
  ``inhabitants'', consider subtyping and supertyping.
\end{itemize}

I don't really see the expression problem in the ``difficult situation'' sense
of the word problem. Instead, I see it in the ``math problem'' sort of way: by
adjusting how you approach things, you can make the most out of what
requirements you need in your design.

\section{Looking Forward}\label{looking-forward}

A lot of frustration in Haskell (and programming in general) lies in trying to
force abstraction and tools to work in a way they weren't meant to. Hopefully
this short run-down can help you avoid going \emph{against} the point of these
design patterns and start making the most of what they can offer!

\section{Signoff}\label{signoff}

Hi, thanks for reading! You can reach me via email at
\href{mailto:justin@jle.im}{\nolinkurl{justin@jle.im}}, or at twitter at
\href{https://twitter.com/mstk}{@mstk}! This post and all others are published
under the \href{https://creativecommons.org/licenses/by-nc-nd/3.0/}{CC-BY-NC-ND
3.0} license. Corrections and edits via pull request are welcome and encouraged
at \href{https://github.com/mstksg/inCode}{the source repository}.

If you feel inclined, or this post was particularly helpful for you, why not
consider \href{https://www.patreon.com/justinle/overview}{supporting me on
Patreon}, or a \href{bitcoin:3D7rmAYgbDnp4gp4rf22THsGt74fNucPDU}{BTC donation}?
:)

\end{document}
